%% annex/p02/first-order-resolution.tex --- proofs for P2.2 resolution

\begin{annexitem}{annex:p2.fol.res.1}{thm:p02-fol-substitution-composition}{Simultaneous-substitution composition}
\begin{proof}
Structural induction proves the term equation.  The homomorphic clauses give
atoms and literals, pointwise images give clauses, and extensionality gives
associativity.
\end{proof}
\end{annexitem}

\begin{annexitem}{annex:p2.fol.res.2}{thm:p02-fol-clause-substitution}{Satisfaction under simultaneous term substitution}
\begin{proof}
Let \(s^{\theta}(x)=\operatorname{val}^{\mathcal M}_{s}(\theta(x))\).
Term induction transfers literal truth under \(s^{\theta}\) to instance truth
under \(s\).  Since \(C\) has a true literal under \(s^{\theta}\), so does
\(C\theta\) under arbitrary \(s\).
\end{proof}
\end{annexitem}

\begin{annexitem}{annex:p2.fol.res.3}{thm:p02-fol-unification}{First-order unification}
\sketchnote{the full lexicographic termination measure and the invariant
calculation for each transformation are omitted.}
Maintain: the original unifiers are precisely \(\theta\eta\), with \(\eta\)
solving the remaining equations.  Every transformation preserves this;
clashes and failed occurs checks admit no finite-term solution.  Elimination,
decomposition, or deletion decreases the lexicographic measure, while
orientation enables one.  Hence termination; at success the invariant gives
an MGU, and at failure no unifier.
\end{annexitem}

\begin{annexitem}{annex:p2.fol.res.4}{thm:p02-fol-resolution-soundness}{Soundness of Resolution and Factoring}
\begin{proof}
Renaming preserves satisfaction using its inverse; substitution does so by
the transformed-assignment lemma.  If an assignment falsified
the resolvent, the parents would force \(A\theta\) and \(\neg B\theta\),
although \(A\theta=B\theta\).  Factoring is analogous.  Induction preserves
models through a derivation; the empty clause is \(\bottom\).
\end{proof}
\end{annexitem}

\begin{annexitem}{annex:p2.fol.res.5}{thm:p02-fol-resolution-completeness}{Refutational completeness of first-order resolution}
\sketchnote{the cited proofs of Herbrand's clause theorem and the lifting
lemma, together with the step-by-step lifted derivation, are omitted.}
Soundness gives the reverse.  For unsatisfiable finite equality-free \(S\),
add a constant if needed.  Herbrand supplies finitely many propositionally
unsatisfiable ground instances; propositional resolution refutes them.
Lifting replaces every ground step by a Resolution or Factoring instance, so
the finite lifted derivation ends in the empty clause.
\end{annexitem}
